% TOP_PROBLEM: {"id": 30006836, "title": "Kaplansky's Idempotent Conjecture", "category_id": 4, "category": {"id": 4, "name": "algebra", "display_name": "Algebra", "description": "Group theory, ring theory, field theory, and algebraic structures.", "slug": "algebra", "order_index": 4, "created_at": "2026-07-31T15:26:25.671Z"}, "status": "open"}
% FIRST_POSTED: 2026-09-25
% ENTRY_KIND: statement_only
% !TeX program = lualatex
% Definition and sources only; this file is not an LLM solution attempt.
\documentclass[11pt]{article}
\usepackage[margin=25mm]{geometry}
\usepackage{fontspec}
\setmainfont{Latin Modern Roman}
\usepackage{amsmath,amssymb,mathtools}
\usepackage{unicode-math}
\setmathfont{Latin Modern Math}
\usepackage{xurl,hyperref}
\hypersetup{colorlinks=true,urlcolor=blue}
\setcounter{secnumdepth}{0}
\setlength{\parindent}{0pt}
\setlength{\parskip}{5pt}
\setlength{\emergencystretch}{3em}
\begin{document}
\section*{\texorpdfstring{Kaplansky's Idempotent Conjecture}{Kaplansky's Idempotent Conjecture}}\label{30006836}
\subsection{Definitions and mathematical statement}
For a field $K$ and group $G$, the ordinary group ring consists of finite sums $\sum_g a_gg$, with multiplication extended bilinearly from $G$. An idempotent is an element $e$ satisfying $e^2=e$. The assertion is
\[
 G\text{ torsion-free},\quad e\in K[G],\quad e^2=e
 \quad\Longrightarrow\quad e=0\text{ or }e=1.
\]
Torsion-free means no nonidentity element has finite order. Fields of every characteristic and groups without finite-generation hypotheses are included. Matrix rings, operator-algebra completions, and infinite formal sums are not in scope. In particular, an idempotent in one of those larger objects is not a counterexample to this scalar finite-support statement.
\par\smallskip\textit{Formulation and concept references:} \href{https://arxiv.org/pdf/1904.04847}{[S1]}.

\subsection{Short English statement}
Are zero and one the only idempotents in the ordinary group ring of every torsion-free group over every field?
\subsection{Sources}
\begin{itemize}
\item[S1]Johan Öinert, Units, Zero-divisors and Idempotents in Rings Graded by Torsion-Free Groups (arXiv:1904.04847v3), Problem 1(c); Gardam (2021), introduction.
\url{https://arxiv.org/pdf/1904.04847}.
\textit{Formulation source.}
\end{itemize}
\end{document}
